\documentclass[11pt]{article}

\usepackage[margin=1.1in]{geometry}
\usepackage{amsmath}
\usepackage{booktabs}
\usepackage{graphicx}
\usepackage{caption}
\usepackage{setspace}
\usepackage{natbib}
\usepackage{xcolor}
\usepackage[colorlinks=true, linkcolor=blue!50!black, citecolor=blue!50!black, urlcolor=blue!50!black]{hyperref}
\usepackage{threeparttable}
\usepackage{float}

\captionsetup{font=small, labelfont=bf}
\setstretch{1.15}

\newcommand{\coef}[1]{\ensuremath{#1}}
\newcommand{\pval}[1]{\ensuremath{(p=#1)}}

\title{\textbf{Do AI Adoption Signals Predict Company Performance?}\\
\large Evidence from a 500-Company Cross-Section of Scores, Filings, and Hiring Data}

\author{
Yixiao Li \and Siru Tao \and Xin Xu \and Hanzhe Hong\\[4pt]
\normalsize Carnegie Mellon University --- AI Capstone Program\\
\normalsize in collaboration with Larridin
}
\date{July 2026 --- Working Draft}

\begin{document}
\maketitle

\begin{abstract}
\noindent
Enterprises are investing heavily in artificial intelligence, yet there is little
systematic evidence on whether observable AI-adoption signals translate into
measurable business outcomes. We study this question using a cross-section of
roughly 500 large U.S.\ public companies, combining (i) proprietary AI-adoption
scores from Larridin's AI Transformation Tracker, (ii) two novel signal families
that we construct from primary sources---LLM-extracted disclosure signals from
10-K filings and an AI-hiring intensity signal built from 30{,}000+ classified
job postings---and (iii) realized financial outcomes from SEC filings and market
data. The analysis sample excludes the five largest AI-semiconductor firms
(Nvidia, Broadcom, AMD, Micron, Intel) so that no handful of names drives the
cross-section; every finding is unchanged on the full sample (appendix).
Three findings emerge. First, the measurement thesis is validated: every score- and
disclosure-based adoption signal in the study is significantly associated with
revenue growth, and the association strengthens with measurement depth. Its
sharpest expression is \emph{narrative concreteness}---a new dimension we
construct within the same evidence-first measurement framework, capturing
whether a firm reports named, deployed use cases with quantified results
rather than aspirational language. Moving from the least to the most concrete
disclosure is associated with 8.0 percentage points higher year-over-year
revenue growth ($p<0.01$), an effect that survives sector, size, and
growth-momentum controls; composite scores, which aggregate several
scale-correlated components, capture the same phenomenon unconditionally,
and isolating their disclosure-concreteness component is what sharpens the
signal. Second, none of the public signals predicts risk-adjusted stock
returns over the subsequent four months, consistent with markets already
impounding public AI-adoption information; margin effects are likewise absent,
indicating that AI currently operates through the growth channel rather than
the cost channel. Third, an AI value-chain decomposition shows the association
is present \emph{within} physical-asset-heavy late adopters
($p<0.05$)---precisely where AI narratives vary most in credibility and
structured measurement is most needed---and that AI-infrastructure suppliers
outperformed sector- and size-matched peers by 32 percentage points over the
sample window. We contribute a validated, reproducible pipeline for measuring
corporate AI adoption from public data, an immediately deployable new scoring
dimension, and the first test--retest reliability estimates for signals of
this class.
\end{abstract}

\section{Introduction}

Every large enterprise now claims an artificial-intelligence strategy. Capital
expenditure guidance, annual reports, and hiring pages are saturated with AI
language, and a measurement industry has emerged to score companies on their
``AI readiness.'' Yet executives, investors, and analysts still lack a rigorous,
data-driven answer to a basic question: \emph{do observable AI-adoption signals
carry information about subsequent business performance, or are they merely
narrative?}

This paper studies that question empirically. Our starting point is the AI
Transformation Tracker built by Larridin, which assigns evidence-backed
1--5 scores to 562 public companies across three pillars---AI adoption,
workforce proficiency, and realized impact---together with an overall maturity
index. We link these scores to market data and SEC fundamentals for a
564-company universe (the Larridin coverage set united with the S\&P 500), and
we extend the signal set with two families of new measures, built at the
sponsor's direction as extensions of the same evidence-first measurement
methodology that underlies the Tracker:

\begin{enumerate}
  \item \textbf{Filing-disclosure signals.} Using a large language model (LLM)
  extraction pipeline with mandatory evidence citation, we score 478 companies'
  most recent 10-K filings on three dimensions: \emph{AI investment intensity}
  (the scale of resources committed to building or buying AI),
  \emph{narrative concreteness} (whether the AI story consists of named,
  deployed use cases with quantified results, versus aspirational boilerplate),
  and \emph{risk-disclosure depth}.
  \item \textbf{AI-hiring intensity.} We collect and individually classify
  30{,}861 job postings across 536 companies in two monthly snapshots
  (June and July 2026), measuring the share of each firm's hiring that targets
  AI-builder roles. The posting classifier is validated against 657
  independently labeled postings (${\sim}90\%$ agreement on the core class),
  and the two snapshots yield the first test--retest reliability estimate we
  are aware of for this signal class (Spearman $\rho = 0.54$).
\end{enumerate}

We then ask three questions. (1) Do AI-adoption signals predict
\emph{operational} outcomes---revenue growth, margin change, productivity? (2)
Do they predict \emph{market} outcomes---risk-adjusted forward returns? (3)
Where in the AI value chain are any effects concentrated?

\paragraph{Findings.} Our headline result is best summarized as
\emph{measurement depth pays}. Every score- and disclosure-based signal in the study---the
composite Tracker scores and our filings dimensions---is
significantly associated with subsequent revenue growth unconditionally
(the hiring signal postdates the outcome quarter and is excluded from this
claim by construction),
validating the underlying thesis that AI adoption is measurable from public
sources and economically meaningful. As controls tighten, the signals then
separate by granularity. Composite scores, which by construction aggregate
several scale-correlated components (deployment breadth, vendor footprint,
enablement infrastructure), largely attenuate as sector and size controls enter. The disclosure-concreteness dimension---the
deepest, most evidence-anchored measurement in the set---survives everything:
its full-specification coefficient of \coef{0.080} \pval{0.009} implies that
moving from the bottom to the top of the concreteness distribution is
associated with 8.0 percentage points higher revenue growth, holding sector,
size, and momentum fixed, and the effect passes a Benjamini--Hochberg false
discovery correction within our pre-specified primary hypothesis family.
Investment intensity is positive throughout but loses significance once size
enters. The practical
reading is constructive: the information is real, and it concentrates in
verifiable specifics---which is precisely what a structured,
evidence-first measurement system is built to capture, and what naive
keyword-counting cannot.

On the market side we find nothing: no signal predicts four-month forward
returns after controls and multiple-testing correction. We interpret this
through the lens of market efficiency---all our signals are constructed from
public information, and the point estimate on the composite score is, if
anything, mildly negative over a window in which early-2026 AI-theme valuations
partially unwound. Margin effects are likewise absent: in this sample, AI
adoption is a top-line story, not (yet) a cost story.

Finally, a four-category AI value-chain classification of the S\&P 500
(infrastructure suppliers, software platforms, data-rich adopters, and
physical-asset-heavy late adopters) reveals two forms of heterogeneity. First,
AI-infrastructure suppliers outperformed sector- and size-matched peers by
approximately 32 percentage points over four months---a striking decomposition
of the H1-2026 AI rally. Second, and more subtly, the concreteness--growth
association is significant \emph{within} the physical-asset-heavy late-adopter
group ($n=214$, $p=0.025$): among companies where AI claims are cheapest to
make and hardest to verify, concrete disclosure best separates genuine adopters
from laggards.

\paragraph{Contributions.} Beyond the empirical findings, we contribute
methodology and infrastructure: (i) an evidence-first LLM extraction protocol
in which every score must cite a verbatim source span (87--90\% of citations
verify against source text), (ii) a validated, repeatable monthly pipeline for
AI-hiring measurement, including classifier validation against independent
labels and cross-month reliability estimates, and (iii) an honest accounting of
what a single-vintage score history can and cannot identify, which we hope is
useful to the emerging AI-transformation-measurement industry.

\section{Data}
\label{sec:data}

\subsection{Universe construction}

Our universe is the union of Larridin's January-2026 coverage set and the S\&P
500. Starting from 567 companies in the Larridin database, we remove test
records and same-company duplicates (e.g., listings under both ``AMD'' and
``Advanced Micro Devices''), yielding 542 distinct companies, of which 14 are
privately held or subsidiaries with no traded equity. We add 22 S\&P 500
constituents not covered by Larridin, for a study universe of 564 companies
spanning all twelve sectors. Every listed company is mapped to its ticker and
SEC Central Index Key through a three-stage procedure (exact name match against
the S\&P constituent list, exact match against the SEC registry, then a
conservative fuzzy pass with manual review); all 542 Larridin companies
resolve. Where one company lists multiple share classes, one line is retained. From
this universe we form the \textbf{analysis sample by excluding the five largest
AI-semiconductor firms}---Nvidia, Broadcom, AMD, Micron, and Intel---whose
extreme H1-2026 revenue and return realizations would otherwise let a few names
dominate a cross-section of this size. All main-text results use this exclusion
sample; the full sample retaining these five firms is reported in Appendix
Table~\ref{tab:fullsample}, where every conclusion is unchanged.

\subsection{Larridin AI-adoption scores}

Larridin's AI Transformation Tracker assigns each company three pillar scores
on a 1--5 scale---\emph{adoption} (deployment breadth, hiring intensity, vendor
footprint, governance), \emph{proficiency} (workforce fluency, enablement
infrastructure, leadership depth), and \emph{impact} (quantified outcomes,
financial linkage, repeatability)---plus a composite \emph{maturity index}.
Scores are generated by an LLM research pipeline over public sources with
per-dimension evidence citations and confidence labels. Our vintage was
assessed on January 18--19, 2026 and covers 562 companies (538 after our
deduplication), with a healthy cross-sectional distribution (maturity mean
3.47, full 1--5 support). Because a single assessment vintage was available at
the time of study, our design is cross-sectional; we return to the implications
in Section~\ref{sec:limitations}.

\subsection{Outcome variables}

\paragraph{Fundamentals.} From SEC XBRL company facts we extract, for each
company, the first fiscal quarter ending after the score date (for most
companies, calendar Q1 2026, reported April--May 2026). Our primary outcome is
\textbf{year-over-year revenue growth} for that quarter (available for 438
companies in the regression sample), which nets out seasonality. Secondary
outcomes are the year-over-year \textbf{operating-margin change} ($n=348$) and
revenue per employee where disclosed. Values are taken as first reported
(avoiding restatement look-ahead) and winsorized at the 1st/99th percentiles.

\paragraph{Market outcomes.} We compute forward returns from the first trading
day after the score date (January 20, 2026) at one- through four-month
horizons, using adjusted closes for 531 tickers. Nine companies delisted through
M\&A or bankruptcy during 2025 are excluded by construction; one mid-window
delisting retains its truncated horizons.

\subsection{Control variables}

Regressions control for \textbf{sector} (twelve fixed effects),
\textbf{firm size} (log market capitalization at the score date, computed from
SEC-reported shares outstanding times the first post-score close for 490
companies, with a current-cap fallback for 36), and \textbf{growth momentum}
(year-over-year revenue growth of the last fiscal quarter ending \emph{before}
the score date, $n=517$)---the observable growth an investor could see at
signal time. The momentum control addresses the central confound that
faster-growing firms both adopt more AI and continue to grow.

\section{Signal Construction}
\label{sec:signals}

\subsection{Filing-disclosure signals}

For each company we retrieve the most recent 10-K from SEC EDGAR, extract
AI-relevant passages by keyword windowing (with section tags so that
risk-factor material is identified as such), and score three dimensions on
anchored 1--5 rubrics using a frontier LLM at temperature zero:

\begin{itemize}
  \item \textbf{AI investment intensity}: the scale of resources committed to
  AI, whether \emph{building} it (AI R\&D, chips, platforms) or
  \emph{buying/deploying} it (compute, cloud, talent). Anchors range from
  ``investment described as minimal or paused'' (1) to ``quantified, sustained,
  or core-business-scale commitment'' (5). Generic R\&D with no stated AI link
  does not count.
  \item \textbf{Narrative concreteness}: 1 = buzzwords only; 3 = named use
  cases without deployment detail; 5 = deployed use cases with quantified
  business results.
  \item \textbf{Risk-disclosure depth}: from no AI risk discussion to
  quantified, governance-backed treatment. A scope guard prevents scoring
  ``absence'' when risk-section material simply was not captured by the
  excerpting step.
\end{itemize}

The protocol is \emph{evidence-first}: a non-null score must cite verbatim
supporting spans, and dimensions without evidence are returned as null rather
than guessed. In an automated audit, 87\% of evidence quotes reproduce verbatim
in the source filings (the remainder being ellipsized quotations). Of 488
companies with a U.S.\ 10-K, 478 yielded scores; ten filings contain no
AI-relevant content, which we record as missing rather than as a low score.
Prompts are versioned; the investment-intensity rubric was revised once
(v1$\to$v2) after review showed the buyer-side framing understated AI
\emph{producers} (e.g., a leading AI-chip designer initially scored~3), and the
full universe was re-scored under v2. Distributionally, concreteness spans the
full 1--5 range; risk-disclosure is heavily concentrated at 4 (75\% of firms),
reflecting the post-2024 standardization of AI risk language---we therefore
treat it as descriptive rather than discriminating.

\subsection{AI-hiring intensity}

We measure AI-directed hiring demand in two monthly snapshots (windows
May~10--June~10 and June~4--July~5, 2026). For each company we retrieve up to
25 recent job-posting search results (60 for the 150 most AI-active companies
in the July snapshot) via a neural web-search API, then classify every result
with an LLM against a four-way taxonomy aligned with Larridin's own labeling
scheme: \emph{AI-builder} (the role's core function is building or operating
AI/ML systems), \emph{AI-user}, \emph{AI-leader}, and \emph{non-AI}, with
conservative defaults, employer-match filtering, and per-posting evidence.
Company-level intensity is the AI-builder share of matched postings
(\emph{builder rate}); 30{,}861 postings were classified in total.

\paragraph{Validation.} Before any full-scale run, the classifier was validated
against 657 postings independently collected and labeled by Larridin's own
crawler-based pipeline: agreement on the AI-builder class is approximately
90\%, and a stronger scoring model yields identical recall on the same inputs,
indicating performance is input- rather than model-limited. At the company
level, our June intensities reproduce the ordering of Larridin's
proof-of-concept measurements for the overlapping companies (Spearman
$\rho=0.83$). Across the two snapshots, company builder rates correlate at
$\rho=0.54$ ($n=226$ companies with at least five matched postings in both
months)---to our knowledge the first published test--retest reliability figure
for a posting-derived AI-hiring signal. Cross-month comparisons use
within-month percentile ranks, as the underlying search index is not
level-stationary across months.

\subsection{AI value-chain classification}

Finally, each S\&P 500 company is assigned one of four AI value-chain
categories---\emph{AI Infrastructure \& Hardware} (36), \emph{AI Software \&
Platforms} (41), \emph{Data-Rich Adopters} (138), and
\emph{Physical-Asset-Heavy / Late Adopters} (285)---by a frontier LLM with a
structured rubric, confidence, and rationale. Unlike the signals above, this
classification draws on the model's general knowledge of each firm rather than
cited primary evidence; we therefore use it only as a segmentation device, not
as a predictive signal.

\section{Methodology}
\label{sec:method}

Our design is a single cross-section: signals measured at or before early 2026
are related to outcomes realized afterward. Signals are transformed to
cross-sectional percentile ranks (robust to outliers and scale), so
coefficients read as the outcome difference between the lowest- and
highest-ranked firm. For each signal $s_i$ and outcome $y_i$ we estimate a
ladder of OLS specifications with HC3 heteroskedasticity-robust standard
errors:
\begin{align*}
\text{(1)}\quad y_i &= \alpha + \beta\, s_i + \varepsilon_i\\
\text{(2)}\quad y_i &= \alpha + \beta\, s_i + \gamma_{\text{sector}(i)} + \varepsilon_i\\
\text{(3)}\quad y_i &= \alpha + \beta\, s_i + \gamma_{\text{sector}(i)} + \delta \log \text{MktCap}_i + \varepsilon_i\\
\text{(4)}\quad y_i &= \alpha + \beta\, s_i + \gamma_{\text{sector}(i)} + \delta \log \text{MktCap}_i + \lambda\, g^{\text{prior}}_i + \varepsilon_i
\end{align*}
where $g^{\text{prior}}_i$ is pre-signal revenue growth. Surviving
specification (4) requires a signal to carry information beyond sector
membership, firm size, and growth persistence. Because a handful of
AI-semiconductor mega-caps posted revenue and return realizations far outside
the rest of the cross-section in this window, our analysis sample excludes the
five largest of them---Nvidia, Broadcom, AMD, Micron, and Intel---so that no
small cluster of names can drive the estimates; the full sample including these
firms is reported as a robustness appendix (Table~\ref{tab:fullsample}) and
leaves every conclusion intact (the headline concreteness coefficient is
\coef{0.080} here versus \coef{0.087} with them). Each estimate uses all firms
with complete data for that specification (missingness is documented and
non-random: banks lack comparable margin data, foreign filers lack 10-Ks); a
complete-case robustness sample is reported in the appendix. Revenue growth was
designated the primary outcome before testing; within the primary family of
five signal tests we apply Benjamini--Hochberg FDR control at $q=0.05$, and we
also report the correction across the full five-signal, three-outcome
(fifteen-test) matrix.
Rank-based information coefficients and quintile portfolio sorts complement the
regressions.

Because job postings are collected live and cannot be reconstructed
retrospectively, hiring signals postdate the fundamental outcome quarter; we
therefore report hiring associations as validation and dynamics rather than
prediction, and we flag that the outcome quarter (typically January--March
2026) overlaps the mid-January score date by roughly two weeks, making
``near-term association'' the precise claim for fundamentals.

\section{Results}
\label{sec:results}

\subsection{Do AI signals relate to operational performance?}

Table~\ref{tab:ladder} reports the specification ladder for the primary
outcome. Three patterns stand out. First, \textbf{every} score- and filings-based
signal is strongly associated with revenue growth unconditionally (column~1):
firms that score high on any AI dimension grew faster. The phenomenon the Tracker is
built to measure is present in the data. Second, the signals separate by
measurement granularity as controls tighten: the composite scores---which
aggregate several components that naturally scale with firm size---remain
significant with sector controls but attenuate once size enters. Third,
\textbf{narrative concreteness, the deepest single dimension in the
measurement stack, survives everything}: its full-specification coefficient of
\coef{0.080} \pval{0.009} implies an 8.0-point revenue-growth gap between the
least and most concrete AI disclosers, holding sector, size, and momentum
fixed. It is the only signal to clear the Benjamini--Hochberg threshold within
the primary family; investment intensity is positive throughout but loses
significance once size enters \pval{0.14}.

\begin{table}[t]
\centering
\begin{threeparttable}
\caption{Specification ladder: AI signals and year-over-year revenue growth}
\label{tab:ladder}
\small
\begin{tabular}{lccccc}
\toprule
 & (1) Univariate & (2) +Sector & (3) +Size & (4) +Momentum & $n_{(4)}$\\
\midrule
Narrative concreteness (ours) & 0.129*** & 0.083*** & 0.075** & \textbf{0.080***} & 395\\
 & (0.000) & (0.010) & (0.018) & (0.009) & \\
Investment intensity (ours) & 0.106*** & 0.063** & 0.048 & 0.041 & 383\\
 & (0.001) & (0.043) & (0.115) & (0.141) & \\
Larridin adoption & 0.108*** & 0.075** & 0.051 & 0.038 & 429\\
 & (0.000) & (0.013) & (0.106) & (0.201) & \\
Larridin maturity index & 0.095*** & 0.058** & 0.032 & 0.024 & 429\\
 & (0.001) & (0.044) & (0.278) & (0.401) & \\
Larridin impact & 0.096*** & 0.062** & 0.042 & 0.025 & 429\\
 & (0.001) & (0.028) & (0.138) & (0.346) & \\
Larridin proficiency & 0.066** & 0.022 & $-0.011$ & $-0.011$ & 429\\
 & (0.015) & (0.493) & (0.749) & (0.735) & \\
AI-hiring builder rate (ours) & 0.026 & 0.021 & $-0.004$ & $-0.027$ & 212\\
 & (0.400) & (0.505) & (0.892) & (0.165) & \\
\bottomrule
\end{tabular}
\begin{tablenotes}\footnotesize
\item Signals in cross-sectional percentile ranks; coefficients are the
outcome difference between lowest- and highest-ranked firms. HC3 robust
$p$-values in parentheses; */**/*** denote $p<0.10/0.05/0.01$. Spec (2) adds
twelve sector fixed effects; (3) adds log market capitalization at the score
date; (4) adds pre-signal year-over-year revenue growth. Hiring postdates the
outcome quarter and is shown for completeness only.
\end{tablenotes}
\end{threeparttable}
\end{table}

The interpretation we place on this pattern is that \emph{specificity is the
signal}. What distinguishes subsequently faster-growing firms, within sector
and size class, is verifiable substance---named systems in production and
quantified results---in their own regulatory disclosures. This is an argument
\emph{for} structured measurement, not against composite scoring: surface
keyword exposure carries no information (risk-disclosure language, now
standardized across firms, discriminates nothing), whereas each step toward
evidence-anchored specificity adds signal. The concreteness dimension was
built inside the same extraction framework as the Tracker's pillars and can be
incorporated into it directly.

\subsection{Is it priced? Market and margin outcomes}

Applying the full specification to four-month forward returns yields no
positive predictability for any signal, and no coefficient in the return
family survives multiple-testing correction. Point estimates are consistent
with a partial unwind of early-2026 AI-theme valuations over this particular
window rather than with any stable relationship.
Margin changes show no association with any signal (all $|t|<1.3$). Taken
together with the revenue results, the evidence points to AI adoption
operating, at this stage, through the \textbf{growth channel}---and to public
adoption information being largely \textbf{impounded in prices} already, as
efficient-markets reasoning would predict for signals built from public
sources.

\subsection{Where in the value chain? Heterogeneity}

Two findings emerge from the value-chain segmentation
(Table~\ref{tab:hetero}). First, the concreteness--growth association is
concentrated where AI credibility is scarcest: it is significant within
\emph{Physical-Asset-Heavy / Late Adopters} (\coef{0.065}, \pval{0.025},
$n=214$) and within \emph{AI Software \& Platforms} (\coef{0.074}, $p=0.049$), while composite scores are insignificant within every category.
Among the mass of ``ordinary'' companies---where AI claims are cheapest to make
and hardest to verify---concrete disclosure best separates genuine adopters
from laggards.

Second, the market leg of the story is a value-chain story: AI-infrastructure
suppliers returned $+44.9\%$ on average over four months, versus roughly flat
for every other category, and the gap survives sector and size controls at
$+31.6$ percentage points ($p=0.005$). Because our analysis sample already
excludes the five largest AI-semiconductors, this premium is carried by the
remaining twenty-five infrastructure names rather than a few marquee stocks
(the full sample including the five gives $+34.2$pp;
Table~\ref{tab:fullsample}). We present it as a descriptive decomposition of
the H1-2026 AI rally---the classification is made with knowledge of these
firms---rather than as an ex-ante trading result; even so, it shows how
narrowly AI's market-value creation concentrated in one segment of the value
chain.

\begin{table}[t]
\centering
\begin{threeparttable}
\caption{Heterogeneity by AI value-chain position}
\label{tab:hetero}
\small
\begin{tabular}{lcccc}
\toprule
 & \multicolumn{2}{c}{Concreteness $\to$ growth} & \multicolumn{2}{c}{Mean 4-month return}\\
\cmidrule(lr){2-3}\cmidrule(lr){4-5}
Category & coef.\ ($p$) & $n$ & raw & $n$\\
\midrule
AI Infrastructure \& Hardware & $+0.149$ (0.305) & 22 & $+44.9\%$ & 25\\
AI Software \& Platforms & $+0.074$ (0.049) & 35 & $-3.8\%$ & 38\\
Data-Rich Adopters & $+0.090$ (0.198) & 101 & $-3.5\%$ & 132\\
Physical-Asset-Heavy / Late & $+0.065$ (0.025) & 214 & $+1.8\%$ & 278\\
\bottomrule
\end{tabular}
\begin{tablenotes}\footnotesize
\item Left panel: within-category regressions of revenue growth on
concreteness rank with size and momentum controls (sector effects omitted
within category). Right panel: mean four-month forward return from the score
date. The infrastructure-vs-late-adopter return gap is $+31.6$pp with sector
and size controls ($p=0.005$); the full sample including the five excluded
semiconductors gives $+34.2$pp.
\end{tablenotes}
\end{threeparttable}
\end{table}

\subsection{Hiring dynamics}

The hiring signal cannot be tested against outcomes that predate it; its
results are methodological and descriptive. Beyond the validation figures of
Section~\ref{sec:signals}, the June--July panel reveals economically sensible
movement: rank-percentile accelerations are led not by technology firms but by
insurers, retailers, and financial-data companies staffing data-science and
AI-engineering roles (verified against the underlying postings), while several
industrial and pipeline firms decelerate. As the monthly panel accumulates,
lead--lag tests of the proposal's central hiring hypothesis---whether AI-hiring
intensity precedes efficiency gains---become feasible.

\subsection{Implications for practitioners}

The results carry three direct implications for AI-transformation measurement.
\emph{First, weight verifiable specificity.} The predictive content of
adoption measurement concentrates in evidence-anchored, deployment-level
detail; incorporating the concreteness dimension into composite scores---or
reporting it alongside them---should sharpen firm-level discrimination,
particularly within peer groups of similar size and sector.
\emph{Second, measure where narratives are cheapest.} The association is
strongest among physical-asset-heavy late adopters, the segment where AI
claims are least verifiable by casual reading and where buyers of measurement
(investors, boards, enterprise customers) face the greatest information
asymmetry. \emph{Third, the panel is the prize.} Every quarterly score vintage
added to a fixed methodology converts cross-sectional associations like ours
into progressively stronger designs---lead--lag tests, within-firm changes,
and eventually out-of-sample validation. The pipelines delivered with this
study (monthly hiring measurement across 536 companies; filings scoring across
478) are built to run on that cadence.

\section{Robustness and Limitations}
\label{sec:limitations}

\paragraph{Robustness.} The headline concreteness result is stable across
winsorization choices, survives sector-neutral rank correlation
(IC $=+0.15$, $p=0.004$; Appendix Table~\ref{tab:ic}), and passes
primary-family FDR control at $q=0.05$; across the broader fifteen-test matrix
it falls just short of the $q=0.10$ threshold---which it cleared only
marginally in the full sample---and we report this plainly. In the strict
complete-case subsample---firms with every signal, outcome, and control
observed, which halves the sample to 196---the coefficient attenuates but stays
positive (\coef{0.061} vs.\ \coef{0.080}) with precision reduced accordingly
($p=0.15$; Appendix Table~\ref{tab:cc}), consistent with a loss of power
rather than sensitivity to sample composition. Score-bin sorts show the
association is monotone (Appendix Table~\ref{tab:sorts}). Five further checks
support the result: a placebo signal (company-name length) run through the
identical pipeline is null ($p=0.84$); in a permutation test, the observed
concreteness coefficient exceeds all but two of 1{,}000 label-shuffled
re-estimates (permutation $p=0.002$); the coefficient remains significant when
any single sector is excluded (worst case: dropping consumer discretionary,
\coef{0.075}, $p=0.029$); it is likewise stable when the five excluded
AI-semiconductors are added back---the full 500-company sample gives
\coef{0.087}, $p=0.007$ (Table~\ref{tab:fullsample}); and re-scoring a random
30-company subsample yields 93\% exact score agreement (100\% within one point;
concreteness rerun rank correlation $0.94$), indicating the LLM scoring itself
is stable. Evidence citation audits,
classifier validation against independent labels, and cross-month reliability
estimates are reported alongside each signal.

\paragraph{Limitations.} Four deserve emphasis. (i) \emph{Single vintage.} One
score snapshot precludes multi-quarter backtesting and walk-forward validation;
our design identifies cross-sectional associations, not time-series
predictability. (ii) \emph{Timing.} The fundamental outcome quarter overlaps
the score date by two weeks, and filing dates straddle it; we therefore claim
near-term association rather than strict prediction, and we restrict return
tests accordingly. (iii) \emph{Disclosure bias.} Concreteness is a disclosure
measure; firms adopting AI quietly are invisible to it, and IR sophistication
may correlate with size in ways our controls only partially absorb.
(iv) \emph{Reverse causality.} Even with momentum controls, faster-growing
firms may write more concrete AI narratives because growth funds deployment;
our estimates should be read as conditional associations. Panel accumulation
across quarterly score vintages---now feasible with the pipelines built
here---is the natural remedy for all four.

\section{Conclusion}

Using a 500-company cross-section that joins proprietary adoption scores, two
newly constructed extension-signal families, and realized outcomes, we find
that AI adoption is already visible in the fundamentals of large U.S.\
companies---and that \emph{how deeply one measures determines how much of it
one sees}. Evidence-based adoption scores identify the phenomenon: highly
rated firms grow faster. Drilling into the measurement stack, the
disclosure-concreteness dimension---concrete, quantified, deployed-use-case
reporting---predicts near-term revenue growth through sector, size, and
momentum controls, margins show no effect yet, and public signals earn no
excess return in a market that appears to have priced them. The value-chain
lens sharpens the picture: infrastructure suppliers captured the market value,
while among late adopters---the majority of the economy---disclosure
concreteness is what separates the firms that are actually growing. The
central lesson is an endorsement of structured, evidence-first measurement
itself: naive AI-mention counting carries no information, composite
evidence-based scores carry the phenomenon, and the deepest evidence-anchored
dimension carries signal that survives every control we impose. The study
also delivers the tooling to compound these findings---a validated,
repeatable measurement pipeline whose quarterly vintages will convert this
cross-section into the panel that the field's causal questions require.

\vspace{1em}
\noindent\textbf{Disclaimer.} This study is exploratory research produced in an
academic--industry collaboration. Nothing herein constitutes investment advice;
backtested or hypothetical results do not represent actual trading.

\section*{Acknowledgments}
We thank Ameya Kanitkar and the Larridin team for data access, API resources,
and weekly guidance, and the CMU AI Capstone faculty for supervision.

\appendix
\renewcommand{\thetable}{A\arabic{table}}
\setcounter{table}{0}

\section*{Appendix: Supplementary Tables}

\begin{table}[H]
\centering
\begin{threeparttable}
\caption{Rank information coefficients (Spearman) by signal and outcome}
\label{tab:ic}
\small
\begin{tabular}{lcccc}
\toprule
 & Revenue growth & Margin change & 4-mo.\ return & Sector-neutral IC\\
Signal & (YoY) & (YoY) & & (revenue growth)\\
\midrule
Narrative concreteness & $+0.222$*** & $+0.019$ & $-0.026$ & $+0.145$***\\
 & (397) & (315) & (456) & \\
Investment intensity & $+0.145$*** & $+0.017$ & $-0.033$ & $+0.051$\\
 & (385) & (306) & (441) & \\
Larridin adoption & $+0.184$*** & $+0.042$ & $-0.073$ & $+0.118$**\\
 & (431) & (342) & (497) & \\
Larridin maturity index & $+0.176$*** & $+0.060$ & $-0.066$ & $+0.104$**\\
 & (431) & (342) & (497) & \\
Larridin impact & $+0.173$*** & $+0.057$ & $-0.049$ & $+0.115$**\\
 & (431) & (342) & (497) & \\
Larridin proficiency & $+0.118$** & $+0.043$ & $-0.074$ & $+0.041$\\
 & (431) & (342) & (497) & \\
AI-hiring builder rate & $+0.012$ & $+0.048$ & $+0.006$ & $-0.070$\\
 & (213) & (174) & (247) & \\
\bottomrule
\end{tabular}
\begin{tablenotes}\footnotesize
\item Spearman rank correlations; observation counts in parentheses.
Sector-neutral IC ranks both signal and outcome within sector before
correlating. */**/*** denote $p<0.10/0.05/0.01$.
\end{tablenotes}
\end{threeparttable}
\end{table}

\begin{table}[H]
\centering
\begin{threeparttable}
\caption{Outcome sorts: revenue growth by signal level}
\label{tab:sorts}
\small
\begin{tabular}{lcc@{\hskip 2em}lcc}
\toprule
\multicolumn{3}{c}{Concreteness score bins} & \multicolumn{3}{c}{Maturity-index quintiles}\\
\cmidrule(r){1-3}\cmidrule(l){4-6}
Score & Mean growth & $n$ & Quintile & Mean growth & $n$\\
\midrule
1 & $-9.4\%$ & 2 & Q1 (low) & $8.5\%$ & 87\\
2 & $7.3\%$ & 100 & Q2 & $8.1\%$ & 86\\
3 & $9.1\%$ & 141 & Q3 & $11.4\%$ & 86\\
4 & $13.4\%$ & 149 & Q4 & $10.0\%$ & 86\\
5 & $46.2\%$ & 5 & Q5 (high) & $16.2\%$ & 86\\
\midrule
\multicolumn{3}{l}{High (4--5) $-$ low (1--2): $+7.5$pp ($t=3.97$)} &
\multicolumn{3}{l}{Q5 $-$ Q1: $+7.8$pp}\\
\bottomrule
\end{tabular}
\begin{tablenotes}\footnotesize
\item Mean year-over-year revenue growth of the first post-score fiscal
quarter, winsorized at 1\%/99\%. Unconditional sorts (no controls); top-quintile
maturity firms grew roughly twice as fast as bottom-quintile firms. Four-month
return sorts show no monotone pattern (Q1 $+2.7\%$ to Q5 $-1.1\%$).
\end{tablenotes}
\end{threeparttable}
\end{table}

\begin{table}[H]
\centering
\begin{threeparttable}
\caption{Complete-case robustness: full specification on the strict subsample}
\label{tab:cc}
\small
\begin{tabular}{lcc}
\toprule
Signal & Complete-case coef.\ ($p$) & Main-sample coef.\ ($p$)\\
\midrule
Narrative concreteness & $+0.061$ (0.148) & $+0.080$ (0.009)\\
Investment intensity & $-0.014$ (0.595) & $+0.041$ (0.141)\\
Larridin adoption & $+0.038$ (0.583) & $+0.038$ (0.201)\\
Larridin maturity index & $+0.026$ (0.666) & $+0.024$ (0.401)\\
AI-hiring builder rate & $-0.027$ (0.224) & $-0.027$ (0.165)\\
\bottomrule
\end{tabular}
\begin{tablenotes}\footnotesize
\item The complete-case subsample ($n=196$) requires every signal, the primary
outcome, and all controls to be observed. Specification (4) of
Table~\ref{tab:ladder}: sector fixed effects, log market capitalization, and
prior revenue growth, HC3 standard errors. Main-sample column is the analysis
sample of Table~\ref{tab:ladder}. The concreteness coefficient attenuates
modestly while precision falls with the halved sample.
\end{tablenotes}
\end{threeparttable}
\end{table}

\begin{table}[H]
\centering
\begin{threeparttable}
\caption{Variable coverage in the analysis sample ($n=513$ listed companies)}
\label{tab:coverage}
\small
\begin{tabular}{lcl}
\toprule
Variable & $n$ & Principal reason for missingness\\
\midrule
Larridin maturity / adoption & 509 & four firms unscored in vintage\\
Narrative concreteness / investment & 457 / 442 & no U.S.\ 10-K (foreign filers); no AI content\\
AI-hiring builder rate & 249 & fewer than five matched postings\\
Revenue growth (YoY) & 434 & non-comparable revenue concepts (banks)\\
Operating-margin change (YoY) & 344 & banks and insurers\\
Four-month forward return & 501 & delistings\\
Log market capitalization & 503 & missing shares data\\
Prior revenue growth & 495 & insufficient reported history\\
\bottomrule
\end{tabular}
\begin{tablenotes}\footnotesize
\item Missingness is non-random and documented; each analysis uses all firms
with complete data for that analysis, with per-cell counts reported throughout.
\end{tablenotes}
\end{threeparttable}
\end{table}

\begin{table}[H]
\centering
\begin{threeparttable}
\caption{Full-sample specification ladder including the five AI-semiconductor firms}
\label{tab:fullsample}
\small
\begin{tabular}{lccccc}
\toprule
 & (1) Univariate & (2) +Sector & (3) +Size & (4) +Momentum & $n_{(4)}$\\
\midrule
Narrative concreteness (ours) & 0.147*** & 0.094*** & 0.084** & \textbf{0.087***} & 399\\
 & (0.000) & (0.005) & (0.011) & (0.007) & \\
Investment intensity (ours) & 0.129*** & 0.076** & 0.057* & 0.048* & 387\\
 & (0.000) & (0.020) & (0.071) & (0.100) & \\
Larridin adoption & 0.115*** & 0.075** & 0.049 & 0.037 & 433\\
 & (0.000) & (0.015) & (0.133) & (0.235) & \\
Larridin maturity index & 0.103*** & 0.060** & 0.032 & 0.024 & 433\\
 & (0.000) & (0.042) & (0.295) & (0.427) & \\
Larridin impact & 0.102*** & 0.063** & 0.041 & 0.023 & 433\\
 & (0.001) & (0.029) & (0.156) & (0.393) & \\
Larridin proficiency & 0.075*** & 0.022 & $-0.013$ & $-0.012$ & 433\\
 & (0.007) & (0.489) & (0.710) & (0.710) & \\
AI-hiring builder rate (ours) & 0.026 & 0.017 & $-0.013$ & $-0.033$ & 215\\
 & (0.393) & (0.590) & (0.712) & (0.108) & \\
\bottomrule
\end{tabular}
\begin{tablenotes}\footnotesize
\item Identical specifications to Table~\ref{tab:ladder}, estimated on the full
sample that \emph{includes} Nvidia, Broadcom, AMD, Micron, and Intel---the five
firms excluded from the main analysis. Every qualitative conclusion is
unchanged: concreteness is the only signal to survive the full specification
(\coef{0.087}, $p=0.007$), and the AI-infrastructure four-month return premium
in this sample is $+34.2$pp ($p=0.003$). HC3 $p$-values in parentheses;
*/**/*** denote $p<0.10/0.05/0.01$.
\end{tablenotes}
\end{threeparttable}
\end{table}

\begin{thebibliography}{9}
\small
\bibitem[Banz(1981)]{banz1981} Banz, R. W. (1981). The relationship between
return and market value of common stocks. \emph{Journal of Financial
Economics}, 9(1), 3--18.
\bibitem[Benjamini and Hochberg(1995)]{bh1995} Benjamini, Y., and Hochberg, Y.
(1995). Controlling the false discovery rate: a practical and powerful approach
to multiple testing. \emph{Journal of the Royal Statistical Society B}, 57(1),
289--300.
\bibitem[Brynjolfsson et al.(2021)]{brynjolfsson2021} Brynjolfsson, E., Rock,
D., and Syverson, C. (2021). The productivity J-curve: How intangibles
complement general purpose technologies. \emph{American Economic Journal:
Macroeconomics}, 13(1), 333--372.
\bibitem[Fama and MacBeth(1973)]{famamacbeth1973} Fama, E. F., and MacBeth,
J. D. (1973). Risk, return, and equilibrium: Empirical tests. \emph{Journal of
Political Economy}, 81(3), 607--636.
\bibitem[Jegadeesh and Titman(1993)]{jt1993} Jegadeesh, N., and Titman, S.
(1993). Returns to buying winners and selling losers: Implications for stock
market efficiency. \emph{Journal of Finance}, 48(1), 65--91.
\end{thebibliography}

\end{document}
